\documentclass[10pt,a4paper]{amsart}
\usepackage[margin=19mm]{geometry}
\usepackage{amsmath,amssymb,mathtools}
\usepackage[scr=rsfs,cal=euler]{mathalfa}
\usepackage{xfrac}
\usepackage[unicode,hidelinks,bookmarks=false]{hyperref}
\newcommand{\ie}{i.e.\ }
\newcommand{\cf}{cf.\ }
\newcommand{\resp}{respectively\ }
\newcommand{\Subsection}{\subsection}
\providecommand{\nobreakdash}{\nobreak\hbox{-}}
\setlength{\emergencystretch}{2em}
\multlinegap0pt
\newtheorem{prop}{Proposition}
\newcommand{\Z}{\mathbb{Z}}
\title{Upper bounds for critical probabilities in \\ Bernoulli Percolation models}
\author{Pablo A. Gomes, Alan Pereira, Remy Sanchis}
\date{Source: arXiv:2106.10388v2 (CC BY 4.0)}
\begin{document}
\maketitle

\noindent\emph{Source and licence.} Upper bounds for critical probabilities in \\ Bernoulli Percolation models. Version arXiv:2106.10388v2 (CC BY 4.0), distributed by arXiv under the Creative Commons Attribution 4.0 International licence, http://creativecommons.org/licenses/by/4.0/, as recorded in the arXiv licence metadata for this exact version and shown on the abstract page https://arxiv.org/abs/2106.10388v2. Full licence notice: \texttt{licenses/arxiv-2106.10388-CC-BY-4.0.txt}.\par\smallskip

\noindent\textit{Editorial scope.} Counted unit: the article's prop prop:triangular (source0.tex lines 348--351). The article's complete original proof of it is delivered with it (source0.tex lines 354--421, one \texttt{proof} environment of 68 lines). No mathematics is written, repaired or re-derived: the statement and the proof are the article's own lines, in the article's own notation, and no retained line is substituted or re-flowed.  No further source-local context is needed: every cross-reference of every retained line resolves inside the two delivered spans.  Nothing was omitted from any retained span, so the package carries no omission record.  No retained line contains a citation, so the wrapper carries no bibliography and no bibliography entry.  No retained line contains a figure environment, an image-inclusion command or drawing code, so no image is delivered and the package declares no asset dependency.  Editorial formatting only, in addition: an AMS \texttt{amsart} preamble that restates the article's own declarations and macros used by the retained lines, namely the environments \texttt{prop} on separate counters, together with the standard packages those lines need.  The retained environments print in this document as Proposition 1 in order of appearance, whereas the article prints the same items with the numbering of its own full document.  The complete native source of the article as distributed in the arXiv e-print for this exact version is bundled unchanged as \texttt{sources/112815/source0.tex}.
\begin{prop}\label{prop:triangular}
Given the parameters $(p_1,p_2,p_3)\in [0,1]$, consider two inhomogeneous bond Bernoulli percolation processes: on the triangular lattice  and on $\Z^3$. Then
\[\theta^b_{\mathbf{T}}(p_1,p_2,p_3)\leq\theta_3(p_1,p_2,p_3).\]
\end{prop}

\begin{proof}
We will construct a dynamic coupling between the percolation process on $\Z^{3}$ with parameters $(p_1,p_2, p_3)$  and an infection process over $\mathbf{T}$.  We will do it in such a way that the law of infected sites in $\mathbf{T}$ is the same as the law of the open cluster of the origin for anisotropic percolation on $\Z^3$ and also that, if the infection process survives, the open cluster of the origin of the process in $\Z^3$ must be infinite. 

Before introduce the coupling, let us set the following auxiliary notation
\begin{equation}\label{eq:tau}
\tau(\pm(1,0))=\pm (1,0,0), \: \tau(\pm(0,1))=\pm (0,1,0) \mbox{ and } \tau(\pm(1,1))=\pm (0,0,1) .
\end{equation}

The coupling will be built based on a susceptible-infected strategy described as follows. First, we declare  the origin of $\mathbf{T}$ as the \textit{initial} infected component. Next, at each time-step, we possibly grow the infected component. %and associate each new vertex $v$ of the infected component in $\mathbf{T}$ to a vertex $x(v)$  in the open cluster of the origin in $\Z^3$.  
Consider a vertex $v$ in the infected component of $\mathbf{T}$ and a neighbor $v+u$ out of the infected component. The vertex $v$ will be associated with some vertex $x(v)$ in the  open cluster of the origin in $\Z^3$. If $\langle x(v), x(v)+\tau(u) \rangle$,  is open, we infect $v+u$  (and write $x(v+u) = x(v) + \tau(u)$).



%If it happens that $u \in \{\pm e_{d}\}$ and $\langle x(v), x(v)+u \rangle$ is closed, we give a second chance to infect $v+u$. For $u=\pm e_d$, in this second chance, we declare $v+u$ infected in $\Z^d$ if $\langle x(v), x(v)\pm e_{d+1} \rangle$ is open in $\Z^{d+1}$. In the last case we set $x(v+u) = x(v)+e_{d+1}$, or $x(v+u) = x(v) - e_{d+1}$, depending on whether $u= e_d$ or $u=-e_d$. 


Formally, the coupling is obtained by the sequence of sets $(I_n, x(I_n), R_n, S_n)_{n \geq 0}$. Here, $I_n$ represents the \textit{infected vertices} in $\mathbf{T}$, $x(I_n)$ represents the vertices in $\Z^{3}$ associated with the infected vertices, and $R_n$ represents the \textit{removed edges} of $\mathbf{T}$. Finally, given $I_n$, $x(I_n)$ and $R_n$, the \textit{susceptible edges set} is given by
\begin{equation*}
    S_{n} := \{ \langle v, u \rangle \in E_{\mathbf{T}}: v \in I_{n} \: \mbox{and} \: u \notin I_{n} \} \cap R_{n}^C.
\end{equation*}
At time $n=0$, we set
\begin{itemize}
    \item $I_0 = \{0\} \subset V_{\mathbf{T}}$; 
    
    \item $R_0 = \emptyset \subset E_{\mathbf{T}}$;
    
    \item $x(0) = 0 \in \Z^{3}$;
    
    \item $S_0 := \{ \langle v, u \rangle : v \in I_{0} \: \mbox{and} \: u \notin I_{0} \} \cap R_{0}^C \:\: = \{ \langle 0, \pm (1,0)\rangle, \langle 0,\pm (1,0)\rangle ,\langle 0,\pm (1,1)\rangle \}$. 
\end{itemize}

This means that, at time $n=0$, only the vertex $0$ is infected, and it can potentially infect any of its neighbors, so all edges containing the origin are susceptible. After that, in each step, an infected vertex tries to infect a non-infected vertex through a susceptible edge (if the latter exists). From now on, we choose an arbitrary, but fixed, ordering of the edges in $\mathbf{T}$. Suppose that $I_n, x(I_n), R_n$ and $S_n$ are already defined. If there is no susceptible edge then the process stops. More specifically, if $S_n = \emptyset$, then for all $k \geq 1$,
    \begin{itemize}
        \item $I_{n+k} = I_n$;
        
        \item $R_{n+k} = R_n$;
        
        \item $S_{n+k} = S_n$.
    \end{itemize}
Otherwise, $S_n \neq \emptyset$. In this case, let $g_n$ be the smallest edge in $S_n$.
Since $g_n \in S_n$, it has to be equal to some $\langle v, v+u_n \rangle$, where $v \in I_n$, $v+u_n \notin I_n$, and $u_n \in \{ \pm (0,1), \pm (1,0), \pm (1,1)\}$. Then, $v$ \textit{infects} $v+u_n$ if $\langle x(v), x(v)+\tau (u_n) \rangle$ is open in $\Z^{3}$ (recall definition of $\tau$ in \eqref{eq:tau}). More precisely, if $\langle x(v), x(v)+\tau (u_n) \rangle$ is open in $\Z^{3}$ then we write
\begin{equation*}
    I_{n+1} := I_n \cup \{v+u_n\},
\end{equation*}
and define
\begin{equation*}
    x(v+u_n) := x(v) + \tau(u_n).
\end{equation*} 
Otherwise, if $\langle x(v), x(v)+ \tau(u_n) \rangle$ is closed in $\Z^{3}$,  we set $I_{n+1}:= I_n$.



In either case, we have explored $g_n$, thus we \textit{remove} it and write
\[ R_{n+1} := R_n \cup \{ g_n \}.\]
Next, to conclude our induction step, we set
\begin{equation*}
  S_{n+1} := \{ \langle v, u \rangle : v \in I_{n+1} \: \mbox{and} \: u \notin I_{n+1} \} \cap R_{n+1}^C.
\end{equation*}
Observe that the function $x: \cup_n I_n \longrightarrow \Z^{3}$ is injective.  
In fact, if $v= (v_1, v_2) \in I_n$, then by construction, we have that $x(v) = (x_1,x_2,x_3)$ satisfies
\[ v_1=x_1+x_3 \mbox{ and }
    v_2=x_2+x_3.\]
Now, observe that the image of $x$ is contained in the  open cluster of the origin $\mathcal{C}^b_0$ of $\Z^{3}$. Since $x$ is injective,  $|\cup_n I_n| \leq |\mathcal{C}^b_0|$.  Also, note that $\cup_{n} I_n$ has the same  law as  $\mathcal{C}^{\mathbf{T}}_0$, where $\mathcal{C}^{\mathbf{T}}_0$ is the  open cluster of the origin in $\mathbf{T}$ with parameters $p_1, p_2,p_3$.  Therefore,
\begin{equation*}
 \theta^b_{\mathbf{T}}(p_1,p_2,p_3) \leq \theta_{3}(p_1,p_2,p_3),
\end{equation*}
and the proof of Proposition \ref{prop:triangular} follows.
\end{proof}

\end{document}
